\scanpage{18}

préschémas \deleted{loc de type fini} \inserted{non ramifiés} séparés sur $S$. Seules les \inserted{deux dernières} conditions de modération se vérifient par des techniques habituelles. \deleted{\unsure{On a|On ne}} \unsure{Celles|Celle} au point générique ne doit pas \unsure{faire|\deleted{faire}} d'ennui, avec le calcul habituel d'obstructions.
% An angular ink line separates the first assertion from the line about the generic point; a wavy stroke runs under the latter.
\margin{\uline{Tout} \uline{\unsure{serait à|suit en|se lit en}} \uline{\unsure{simplifier|simplifiant}}}
% slanted

Reste le cas où $S$ est local noeth irréductible de dim 1, avec Idéal $J$ annulé par le nilradical, d'où $S_0=\mathrm V(J)$, et où on dispose de $u_0:X_0\to Y_0$ (tel que $u_0\varphi_0=v_0$) et un prolongement $u_{\unsure{y|\eta}}:\deleted{\unsure{X}}X_y\to Y_y$ ($y$ le pt générique de $S$) tel que $u_y\varphi_y=v_y$, $u_{y0}=u_{0y}$. Question de prolonger $u_0$ en $u$ tel que $u\varphi=v$, \deleted{\unsure{u}} Soient $P(u_0,S),P(u_0,S_y);P(v_0,S),P(v_0,S_y)$ les prolongements de $u_0$ resp. $v_0$ \deleted{sur} $S$, \deleted{\unsure{resp|et}}\inserted{et} de $u_{0y}$ resp. $v_{0y}$ sur $S_y$, d'où le diagramme commu
% Vertical marginal rule beside the question about prolongation.
\diagram{$(1)$\quad\begin{tikzcd}[ampersand replacement=\&]
 P(u_0,S)\arrow[r]\arrow[d]\&P(u_0,S_y)\ {*}\arrow[d]\\
 {*} \ P(\deleted{\unsure{u}}v_0,S)\arrow[r,dash]\& P(v_0,S_y)
\end{tikzcd}}

d'ens \unsure{form.|form} princ. homogènes, sous les groupes
\diagram{$(1')$\quad\begin{tikzcd}[ampersand replacement=\&]
 \mathrm{Hom}(M^{\deleted{\unsure{v|*}}}_{u_0},J)\arrow[r]\arrow[d]\&\mathrm{Hom}(M_{u_0},J_y)\arrow[d]\\
 \mathrm{Hom}(M_{v_0},J)\arrow[r]\&\mathrm{Hom}(M_{v_0},J_y)
\end{tikzcd}}

Par hypothèse on tient des éléments aux coins
% archive sheet number: 17
